%% kapa: exam=1 round=32 year=2021 period=1 form=A subject=real_estate subject_ko=부동산학원론 seq=38 number=118
%% source: Q-Net 감정평가사 기출문제 (cst00309) artlSeq=5208397 file=extracted/1008/5208397/A형/2021년도 제32회 감정평가사 1차 1교시 A형.pdf
%% conversion: kordoc 4.24.51 -> tools/gen1.py (auto); review=proofread-sonnet; flags=none
\begin{question}{118}
다음 자료를 활용하여 공시지가기준법으로 평가한 대상토지의 단위면적당 가액은? (단, 주어진 조건에 한함)
\begin{framedtext}
○ 대상토지 현황: A시 B구 C동 175번지, 일반상업지역, 상업나지 \par
○ 기준시점: 2021.04.24 \par
○ 비교표준지: A시 B구 C동 183번지, 일반상업지역, 상업용 \par
\hspace*{4em}2021.01.01.기준 공시지가 6,000,000원/m$^2$ \par
○ 지가변동률(2021.01.01\textasciitilde 2021.04.24): A시 B구 상업지역 2\% 상승함 \par
○ 지역요인: 비교표준지와 대상토지는 인근지역에 위치하여 지역요인 동일함 \par
○ 개별요인: 대상토지는 비교표준지에 비해 가로조건에서 5\% 우세하고, 환경조건에서 10\% 열세하며, 다른 조건은 동일함(상승식으로 계산할 것) \par
○ 그 밖의 요인 보정: 대상토지 인근지역의 가치형성요인이 유사한 정상적인 거래사례 및 평가사례 등을 고려하여 그 밖의 요인으로 50\% 증액 보정함 \par
\end{framedtext}
\begin{choices}
  \item 5,700,000원/m$^2$
  \item 5,783,400원/m$^2$
  \item 8,505,000원/m$^2$
  \item 8,675,100원/m$^2$
  \item 8,721,000원/m$^2$
\end{choices}
\end{question}
\answer{4}
